\section{Worldviews}

\begin{quotex}
Metaphysical mutations --- that is to say radical, global transformations in the values to which the majority subscribe --- are rare in the history of humanity. The rise of Christianity might be cited as an example. Once a metaphysical mutation has arisen, it tends to move inexorably toward its logical conclusion. Heedlessly, it sweeps away economic and political systems, aesthetic judgments and social hierarchies. No human agency can halt its progress --- nothing except another metaphysical mutation. It is a fallacy that such metaphysical mutations gain ground only in weakened or declining societies. When Christianity appeared, the Roman Empire was at the height of its powers: supremely organized, it dominated the known world; its technological and military prowess had no rival. Nevertheless, it had no chance. When modern science appeared, medieval Christianity was a complete, comprehensive system which explained both man and the universe; it was the basis for government, the inspiration for knowledge and art, the arbiter of war as of peace and the power behind the production and distribution or wealth --- none of which was sufficient to prevent its downfall.
\flright{\textsc{Michel Houellebecq}, \emph{Elementary Particles}}
\end{quotex}


What Mr. Houellebecq does not mention is that we are in the process of another ``metaphysical mutation'': the post-modern globalist worldview. The process whereby one worldview is overtaken by another is unclear. In an earlier post (Evola on Christianity in Europe\footnote{\url{https://gornahoor.net/?p=27}}), Julius Evola attributed the rise of Christianity during the Roman Empire to a "miracle". One can idly hope for another miracle or else participate in the process of restoring a Traditional worldview.

\paragraph{Addendum 2020}


Perhaps Houellebecq addressed the issue in \textit{Submission}\footnote{\url{https://gornahoor.net/?p=8417}}.

The Miracle that Evola referred to was Constantine\footnote{\url{https://gornahoor.net/?p=11425}}`s vision of Christ before his big battle.

The futility of the human race to control events is the reason for a miracle, not a catastrophe as some desire. So what do you do while you are waiting?

\url{https://www.youtube.com/watch?v=lL2DzqrUNQc} 

\flrightit{Posted on 2008-04-14 by Cologero}

\begin{center}* * *\end{center}

\begin{footnotesize}
\begin{sffamily}
\texttt{Michael on 2020-04-15 at 10:22 said:}

``The futility of the human race to control events is the reason for a miracle, not a catastrophe as some desire. So what do you do while you are waiting?''

Waiting for that, we work on ourselves according to what is given to us, but in that process we start to see more. While the human race's attempts to control events is futile, history is directed through conscious agents. So while we wait, we learn about the forces, agents, and our place among it all? Perhaps then we remember what we were supposed to do.

\end{sffamily}
\end{footnotesize}

